% ============================================================================
% SafeBARS — Multimodal Ethical Mirror — CHI 2027 draft outline + Related Work
% Format: ACM acmart (CHI template). Drop into Overleaf; verify track/length with PCS.
% TODO before submission: verify all author names + the SafeBARS'23 citation on ACM DL.
% This version reframes the paper around the teacher's explicit requirement:
%   (1) MULTIMODAL interaction (Taheri et al. CHI'26 named text-only as a limitation),
%   (2) USER SELF-DISCOVERY of blind spots (not AI telling them),
%   (3) evidence the app CHANGES THE USER'S MIND (not just generates a report).
% ANCHORED ON THE LIVE DEPLOYED SYSTEM: https://safebars-ai.pages.dev/safebars/mirror
%   (Flask + LLM; 9 ethics lenses w/ coverage states; bounded role-probe agents
%    [Perspectra-like]; scenarios; ethics-guided tension queue; deterministic
%    deep-audit). The multimodal self-discovery layer is ported from the prototype
%    into this live site (stakeholder map / value-conflict heatmap / scenario branch
%    tree + "did you anticipate this group?" + before/after re-render).
% DESIGN REFINEMENTS INFORMED BY 4 CHI'26 PAPERS (Vibe Check / RECALLbot /
%    Hear You in Silence / Perspectra): user-selectable expert lenses, agentic
%    cross-session memory, reciprocal disclosure, context-aware pacing, and an
%    explicit critical-vs-sycophantic alignment manipulation. See \S System + RQ6.
% TODO before submission: verify all author names + the SafeBARS'23 citation on ACM DL
%    (ACM Open Access landing blocked full-text fetch; VERIFY against DL).
% ============================================================================
\documentclass[manuscript,screen]{acmart}

\title{SafeBARS Re-Mirrored: A Multimodal Ethical Mirror That Lets Researchers
Discover Their Own Blind Spots}
% TODO: finalize title; add authors / affiliations

\begin{abstract}
% TODO (fill after study): SafeBARS (CHI'23) introduced the "ethical mirror" that
% surfaces tensions between a system's design and ethical considerations. We extend it
% into a *multimodal* mirror in which researchers confront the consequences of their
% OWN design through interactive visualization and scenario simulation. Crucially, the
% mirror withholds prescription: it renders gaps (a stakeholder the researcher never
% named, a value the design leaves uncovered) so the realization is the user's, not the
% AI's. Across a between-subjects study (multimodal vs. text-only mirror) we show that
% the multimodal mirror produces more self-discovered blind spots and measurable
% mind-change -- operationalized as pre/post affected-group listing plus *behavioral*
% plan revision, not a generated report. We contrast this with recent work in which
% AI dialogue shifts bias recognition but remains text-only and where self-reflection
% emerges only by chance.
\end{abstract}

\begin{document}
\maketitle

% ---------------------------------------------------------------------------
\section{Introduction}
% OUTLINE (bullets to expand into prose):
% - Ethical blind spots in research design are hard to see *because they are the
%   researcher's own*: the designer rarely notices who their system affects.
% - SafeBARS (CHI'23) = pre-LLM mirror; reflects only what is written; text-driven.
% - Recent CHI'26 work shows AI dialogue can shift bias recognition (Taheri et al.),
%   but that work is (a) TEXT-ONLY and names multimodality as its own key limitation,
%   (b) about recognizing bias in *other people's* social interactions, and
%   (c) found self-reflection only as a serendipitous side-effect -- it was not designed for.
% - The teacher's requirement, and our contribution: make the mirror MULTIMODAL,
%   make self-discovery DELIBERATE, and PROVE it changes the user's mind.
% - We contribute: (1) a multimodal ethical mirror, (2) a designed self-discovery
%   mechanism ("critical friction" manufactured multimodally), (3) a closed loop with
%   behavioral mind-change evidence, (4) the LLM engine + deterministic deep-audit behind it.
% - NOTE for Erick meeting: tool is deployed & working; this paper closes the loop with a study.

% ---------------------------------------------------------------------------
\section{Related Work}
Recent CHI research has rapidly advanced LLM-based conversational agents across four
dimensions directly relevant to our work. \citep{vibecheck26} show that an agent's
personality expression follows an inverted-U; \citep{recallbot26} demonstrate that
agentic memory combined with reciprocal self-disclosure deepens human--chatbot
relationships; \citep{hearinsilence26} find that context-aware pacing promotes
self-disclosure and affective trust; and \citep{perspectra26} show that letting users
choose and invite expert agents enhances critical thinking. These papers share one
goal: optimizing the agent to build rapport, trust, and intimacy. We borrow their
\emph{mechanisms} (user-selectable perspectives, transparent memory, paced guidance)
while inverting the objective: a critical reflection tool must resist sycophancy and
preserve critical distance rather than maximize liking.

\paragraph{AI-mediated dialogue and bias recognition.}
Most directly adjacent, \citep{taheri26} (CHI'26) show that AI-mediated \emph{dialogue}
shifts recognition of ableist bias more than a text-only reading control, and that
resisting a biased AI coach can forge sharper moral boundaries -- what they call
\emph{critical friction}. Three facts make their work both our anchor and our point of
departure. \emph{First}, their intervention is \textbf{text-only}: they explicitly state
as a limitation that their "text-based dialogue \dots lacks the multimodal richness of
face-to-face interaction" and call for "virtual reality or voice-based agents" as future
work \cite{taheri26}. \emph{Second}, their self-reflection finding was \textbf{serendipitous}
-- an "unintended" learning mechanism, not a designed one. \emph{Third}, they measure
recognition of bias in \emph{others'} interactions via vignette ratings; they do not
measure whether the tool changes the participant's \emph{own} design, nor capture
behavioral revision. Our work takes their strongest mechanism (critical friction) and
\emph{designs for it multimodally}, on the researcher's \emph{own} artifact, with
\emph{behavioral} mind-change as the outcome.

\paragraph{SafeBARS and the ethical mirror.}
We build on the CHI'23 SafeBARS ethical mirror \citep{safebars23}, a pre-LLM tool that
surfaces tensions between a design and ethical considerations and asks researchers to
respond. We inject LLMs (multi-role ethics probes, a deterministic deep-audit layer that
catches risks researchers did not write down) and -- the contribution of this paper --
turn the mirror from a \emph{text report} into a \emph{multimodal, self-discovery loop}.
To our knowledge no reflection tool has (a) used multimodal rendering to trigger
user-owned realization of ethical blind spots, nor (b) measured behavioral mind-change
rather than generated a report.

% ---------------------------------------------------------------------------
\section{SafeBARS: The Multimodal Ethical Mirror}
% Engine (unchanged from prior work): Flask + LLM providers; multi-role ethics probes;
% deterministic deep-audit. What is NEW for CHI'27 is the front-end interaction loop.
\subsection{Engine (background)}
Flask backend + pluggable LLM providers. The system is \emph{deployed and publicly
accessible} at \texttt{safebars-ai.pages.dev/safebars/mirror}. Multi-role ethics probes generate tensions
across ethics lenses (privacy, fairness, autonomy, \dots); researchers choose which
"expert" lenses to enable (mirrors \citep{perspectra26}). A deterministic deep-audit
layer (offline, no LLM budget) catches contradictions, high-risk domains, and
unwritten risks. The engine outputs structured tensions, lens coverage states, bounded
role probes, and suggested scenarios -- the raw material the multimodal layer renders.

\subsection{The Multimodal Layer (contribution)}
Instead of presenting tensions as a text list or a generated PDF, the mirror renders
three coordinated views from the engine output:
\begin{itemize}
  \item \textbf{① Stakeholder map} (interactive SVG): the researcher's design at the
    center; green nodes are groups the researcher named; \textbf{red, pulsing nodes are
    groups the engine inferred the researcher did \emph{not} name}. The gap is seen,
    not stated.
  \item \textbf{② Value-conflict heatmap}: ethics lenses (rows) $\times$ coverage state
    (Missing / Claimed / Reasoned / Action-linked), colored. Red rows = values the
    design leaves uncovered. This is evidence-coverage, explicitly \emph{not} an ethics
    score.
  \item \textbf{③ Scenario branch tree}: a design choice branches into "as designed
    $\to$ consequence for a group" vs. "with a safeguard $\to$ mitigated". The user
    sees the fork their decision creates.
\end{itemize}

\subsection{Self-Discovery, by Design (not by accident)}
Taheri et al.\ found self-reflection only when participants happened to resist a biased
coach. We \emph{manufacture critical friction}: the mirror highlights one red
stakeholder node and \textbf{asks} -- "Did you anticipate this group?" -- rather than
declaring a problem. The user must actively predict / acknowledge; only then does the
mirror reveal the consequence. The realization is the user's; the AI withholds the
verdict. This is the core difference from a report-generating app.

\subsection{Co-Revision With AI Help (not auto-fix)}
When the user responds, the AI offers \textbf{scaffolding}, not a fix
\citep{taheri26}: multiple resolution options (revise / contest-with-evidence /
consult-stakeholder / retain-with-rationale) the user can adapt. After the user writes
the revision, the mirror \textbf{re-renders} the stakeholder map and heatmap
(before $\to$ after), closing the loop and making the change visible and owned.

\subsection{Closed Loop}
Input design $\to$ multimodal mirror $\to$ user self-discovery $\to$ co-revision with
AI scaffold $\to$ re-mirror (before/after). Mind-change is measured at the seam
(pre/post affected-group listing + the actual revised plan as behavioral evidence).

\subsection{Agentic Refinements Informed by CHI 2026}
The live mirror already embodies one CHI'26 mechanism -- \emph{bounded multi-role
probes}~\citep{perspectra26} (researchers are shown several ethics standpoints). We
extend it with four refinements drawn from the CHI'26 agentic-AI literature, each
borrowing a \emph{mechanism} while inverting the field's rapport-optimizing goal:
\begin{itemize}
  \item \textbf{User-selectable experts (Perspectra).} Researchers explicitly
    \emph{invite} the expert lenses/roles they want challenged by, rather than
    receiving a fixed set. Hypothesis: user-chosen perspectives deepen critical
    engagement (tested as a within-condition probe).
  \item \textbf{Agentic memory (RECALLbot).} The mirror remembers a researcher's
    prior commitments and revisions across sessions and can recall ``last time you
    committed X but did not act on it'' -- continuity of ethical attention.
  \item \textbf{Reciprocal disclosure (RECALLbot).} Beyond the boundary notices
    already in the UI, the mirror states its own limits up front: ``I am an AI; I
    infer stakeholders you may not have considered; verify them.'' Calibrates trust.
  \item \textbf{Context-aware pacing (Hear You in Silence).} After a self-discovery
    moment or a hard tension, the mirror pauses rather than surfacing the next
    tension immediately; it batches tensions by the user's revealed cognitive load.
  \item \textbf{Alignment manipulation (Vibe Check).} The mirror's persona is
    deliberately \emph{critical, non-sycophantic}. We expose this as a study
    condition -- \emph{critical} vs.\ \emph{sycophantic/agreeable} alignment (RQ6) --
    to test whether non-sycophancy trades felt agreement for greater critical
    distance and behavioral mind-change, inverting Vibe Check's rapport agenda.
\end{itemize}

% ---------------------------------------------------------------------------
\section{Research Questions}
% Modeled on Taheri's "does X improve Y" structure but push past text-only + report.
\begin{itemize}
  \item \textbf{RQ1 (Multimodal > text-only, cf. \citep{taheri26}):} Does a multimodal
    ethical mirror produce more \emph{self-discovered} blind spots and greater
    mind-change than a Taheri-style text-only dialogue mirror? (Primary outcome:
    self-discovery utterances + behavioral plan revision.)
  \item \textbf{RQ2 (Self-discovery mechanism, cf. \citep{taheri26} critical friction):}
    When the system surfaces a gap \emph{multimodally and without prescribing the
    problem}, do users report more "I didn't realize" moments and show greater
    \emph{ownership} of revisions than when the system states the problem directly?
    (Prescription vs no-prescription within-subject contrast.)
  \item \textbf{RQ3 (Co-revision, not auto-fix):} Does user-led revision with AI
    scaffolding (vs AI auto-fixing the plan) yield higher-quality, more durable
    revisions and retained agency? (Revision-quality rubric + agency self-report.)
  \item \textbf{RQ4 (Does it change the mind, not just report?):} Does the loop change
    researchers' actual coverage -- pre/post affected-group listing \emph{and}
    behavioral plan change -- rather than only generating a document? (This is the
    teacher's "must change the user's mind" bar.)
  \item \textbf{RQ5 (engine characterization, reuse \citep{perspectra26} multi-role idea):}
    Does engine model choice affect which tensions surface? (Answered by the
    computational simulation in \S\ref{sec:sim}; secondary, not the human-outcome claim.)
  \item \textbf{RQ6 (alignment \& critical distance, cf. \citep{vibecheck26}):} Does a
    deliberately \emph{non-sycophantic} mirror (vs.\ a sycophantic one) lower felt
    agreement yet raise critical distance and behavioral mind-change? (Manipulated as
    the alignment condition; probes the inverse of Vibe Check's rapport-optimizing goal.)
\end{itemize}

% ---------------------------------------------------------------------------
\section{Method}
\subsection{Study 1 -- Human user study (primary, answers RQ1--RQ4)}
\noindent\textbf{Design.} Primary between-subjects: \emph{Multimodal mirror} (deployed at
\texttt{safebars-ai.pages.dev/safebars/mirror}, with the three-view self-discovery
layer) vs.\ \emph{Text-only mirror} (Taheri-style: same engine output as a text list /
chat, no visualization). A 2$\times$2 crosses Modality with \emph{Alignment}
(\emph{critical/non-sycophantic} vs.\ \emph{sycophantic/agreeable}) to answer RQ6. A
within-subject prescription contrast (mirror states the problem vs.\ withholds it)
addresses RQ2. Pre $\to$ intervention $\to$ post.

\noindent\textbf{Participants.} Friends pilot (8--12, per SURF guidance) for
feasibility, then a Prolific sample sized for the between-subjects effect (target
$\ge$60--80, matched to Taheri's 160 if resources allow). Eligibility: researchers /
grad students who have designed or are designing an AI/HCI system.

\noindent\textbf{Procedure.} (1)~\emph{Pre}: participant lists the groups their design
affects and rates its ethical robustness. (2)~\emph{Intervention}: submit a real or
provided plan; experience the assigned mirror condition; for RQ2, half receive
prescription, half withhold. (3)~\emph{Self-discovery probe}: the mirror asks "did you
anticipate this group?"; we log "I didn't realize" utterances. (4)~\emph{Co-revision}:
participant revises with AI scaffolding. (5)~\emph{Post}: re-list affected groups; we
code the \emph{actual revised plan} for whether the missed stakeholder / safeguard was
added (behavioral mind-change).

\noindent\textbf{Measures.} (1)~Self-discovered blind spots (coded "I didn't realize"
vs "the system told me"). (2)~Mind-change: $\Delta$ affected-group count (pre/post) +
behavioral revision count. (3)~Revision quality rubric (independent raters).
(4)~Agency / ownership self-report. All coding schemes released with the paper.

\noindent\textbf{Analysis.} Between-subjects ANOVA / mixed models on the above; code
self-discovery vs AI-prescribed utterances; report effect sizes. RQ1--RQ4 answered here.

\noindent\textbf{Instruments (released).} The full protocol -- consent form, step-by-step
script, the 7-construct post-questionnaire (self-discovery, mind-change, agency, critical
distance, trust, aha, manipulation checks) with Likert items, the 2-coder coding scheme
(self-discovery attribution, mind-change, revision-quality rubric; target IRR $\ge$.70), the
analysis plan, and the pilot runbook -- is specified in \texttt{paper/study1\_protocol.md}
(v1.0) with a participant-facing \texttt{.docx} packet. Pre-registration on OSF covers
H1/H3/H4 and the primary DVs.

\subsection{Study 2 -- Computational simulation (secondary, characterizes the engine)}
\label{sec:sim}
\noindent\textbf{Conditions.} For each of six corpus plans we run the mirror engine with
different LLM variants as the \emph{engine}, then prompt the same variant \emph{as a
simulated researcher} to respond to every tension. This characterizes how engine choice
shapes surfaced tensions and critical-distance behavior (RQ5), and probes the engine
without recruiting participants. It is a system characterization, \emph{not} a
human-outcome claim; scripts released (\texttt{modules/multi\_model\_eval.py}).

\subsection{Why a Simulation, and Its Limits}
\label{sec:whysim}
The simulation does \emph{not} substitute for Study 1. LLM-simulated researchers
approximate but do not reproduce human moral reasoning, surprise, or accountability.
We position it as a design probe and engine characterization only; all human-outcome
claims come from Study 1.

% ---------------------------------------------------------------------------
\section{Expected Contributions}
\begin{itemize}
  \item (1) A \textbf{multimodal ethical mirror} (interactive visualization +
    scenario simulation) -- going beyond Taheri's text-only dialogue, which they named
    as their own key limitation.
  \item (2) A \textbf{designed self-discovery mechanism}: critical friction manufactured
    multimodally, so users own the realization instead of being told (vs.\ Taheri's
    serendipitous self-reflection).
  \item (3) \textbf{Behavioral mind-change evidence}: the loop is shown to change
    researchers' actual design coverage, not merely to emit a report.
  \item (4) The LLM multi-role engine + deterministic deep-audit behind the mirror
    (non-sycophantic by design, inverting the rapport-optimizing agenda of
    \citep{vibecheck26,recallbot26,hearinsilence26,perspectra26}).
\end{itemize}

% ---------------------------------------------------------------------------
\section{Limitations and Future Work}
\label{sec:limits}
\begin{itemize}
  \item Study 1's friends pilot is small; the Prolific sample is needed before any
    general claim. The simulation (Study 2) is a system characterization, not evidence
    about humans (see \S\ref{sec:whysim}).
  \item Like Taheri et al., a single session cannot establish \emph{durability} of
    mind-change; longitudinal follow-up is future work.
  \item Voice / VR modalities (Taheri's suggested future environments) are sketched but
    not yet implemented; we prioritize interactive visualization as the first modality.
  \item LLM-as-judge (Study 2) validity: we disclose the rubric and use a deterministic
    fallback; human ratings triangulate the engagement scores.
\end{itemize}

% ---------------------------------------------------------------------------
\section{Conclusion}
% TODO after results.

% ============================================================================
\begin{thebibliography}{9}
% VERIFY all author names + SafeBARS'23 citation against ACM DL before submission.
\bibitem[SafeBARS(2023)]{safebars23}
TODO: SafeBARS CHI'23 authors.
\newblock SafeBARS: ... (ethical mirror).
\newblock In \emph{Proc. CHI '23}.
\newblock \url{https://doi.org/VERIFY}

\bibitem[Taheri et al.(2026)]{taheri26}
A.~Taheri, H.~El Alaoui, P.~Carrington, and J.~P.~Bigham.
\newblock "I followed what felt right, not what I was told": Autonomy, Coaching, and
Recognizing Bias Through AI-Mediated Dialogue.
\newblock In \emph{Proc. CHI '26}, pp.\ 1--23.
\newblock \url{https://doi.org/10.1145/3772318.3791078}

\bibitem[Rahman and Desai(2026)]{vibecheck26}
S. Rahman and S. Desai.
\newblock Vibe Check: Understanding the Effects of LLM-Based Conversational Agents'
Personality and Alignment on User Perceptions in Goal-Oriented Tasks.
\newblock In \emph{Proc. CHI '26}.
\newblock \url{https://doi.org/10.1145/3772318.3790388}

\bibitem[Jiang et al.(2026)]{recallbot26}
Y. Jiang et al.
\newblock RECALLbot: Designing Agentic Memory and Reciprocal Disclosure for
Human-Chatbot Relationships.
\newblock In \emph{Proc. CHI '26}.
\newblock \url{https://doi.org/10.1145/3772318.3790714}

\bibitem[Jiang et al.(2026)]{hearinsilence26}
X. Jiang et al.
\newblock Hear You in Silence: Designing for Active Listening in Human Interaction
with Conversational Agents Using Context-Aware Pacing.
\newblock In \emph{Proc. CHI '26}.
\newblock \url{https://doi.org/10.1145/3772318.3791238}

\bibitem[Liu et al.(2026)]{perspectra26}
Y. Liu et al.
\newblock Perspectra: Choosing Your Experts Enhances Critical Thinking in Multi-Agent
Research Ideation.
\newblock In \emph{Proc. CHI '26}.
\newblock \url{https://doi.org/10.1145/3772318.3791560}
\end{thebibliography}

\end{document}
